% NeurIPS-style paper for Cortex v1
% Based on NeurIPS 2025 formatting guidelines
\documentclass{article}

% --- NeurIPS core formatting ---
\usepackage[final]{neurips_2025}  % [final] for camera-ready; remove for anonymous submission

% --- Standard packages (NeurIPS-approved) ---
\usepackage[utf8]{inputenc}
\usepackage[T1]{fontenc}
\usepackage{hyperref}
\usepackage{url}
\usepackage{booktabs}
\usepackage{amsfonts}
\usepackage{amsmath,amssymb}
\usepackage{nicefrac}
\usepackage{microtype}
\usepackage{xcolor}
\usepackage{algorithm}
\usepackage{algpseudocode}
\usepackage{enumitem}

% --------------------------------------------------------------------------
\title{Cortex v1: A Geometric Controller Coupled to a\\Matrix Product State Information Filter}

\author{%
  Chetan Patil\\
  Independent Researcher\\
  \texttt{chetanxpatil@github}
}

\begin{document}

\maketitle

% =========================================================================
\begin{abstract}
We describe Cortex~v1, a system that couples a deterministic geometric controller (a $3\!\times\!3\!\times\!3$ cubic lattice with 24 rotation symmetries) to a stochastic information filter (a Matrix Product State quantum simulator governed by an entropy-aware polarity mechanism).
Text is processed word-by-word: each word maps to a rotation via a semantic bridge (GloVe embeddings projected into SO(3)), producing a scalar signal~$\alpha$ that modulates the MPS bond-truncation threshold.
We report measured results on three axes:
(i)~MPS simulator scaling from 5 to 1\,000 qubits for structured entanglement,
(ii)~governor-controlled circuits maintaining zero truncation error at 1\,000 qubits, and
(iii)~a retrieval benchmark where $\alpha$-triage achieves 100\% fact recall versus 30\% for FIFO and LRU under identical memory constraints, and
(iv)~a TFIM quantum quench benchmark where the polarity governor reduces mean absolute error in $\langle Z_i \rangle$ by 43.2\% relative to a fixed entropy ceiling under identical $\chi_{\max}$ budget, and
(v)~a Ramsey R(3,3) geometric sandbox where an emergent nonlinear tension field on the $3\!\times\!3\!\times\!3$ lattice detects all $K_3$ violations across 32,768 colorings of $K_6$ (100\%) and a self-healing Simulated Annealing loop finds valid $K_5$ colorings on 100\% of starts while correctly flooring above $T=0$ for $K_6$ (consistent with $R(3,3)=6$), and
(vi)~an R(5,5) heuristic benchmark where a triangle adjacency kernel reduces monochromatic $K_5$ counts by 34\% relative to a cubic density kernel on $K_{43}$ and $K_{44}$ (best: 484/962,598 and 573/1,086,008 respectively); no certificate is claimed.
We state explicitly what the system does and does not do, and provide all code and test suites for independent verification.
\end{abstract}

% =========================================================================
\section{Introduction}

This paper describes a system, not a claim.  The system has three components: a cubic lattice that produces a geometric signal from rotations, a quantum simulator that processes information under an entropy budget, and a bridge that connects the two.  We report what was measured.  Interpretation is left to the reader.

The motivating question is whether a geometric signal derived from word embeddings can control an information-theoretic filter in a way that outperforms content-blind baselines.  The benchmark in Section~\ref{sec:benchmark} provides one data point; the scaling results in Section~\ref{sec:scaling} provide another.  Neither constitutes a general proof.

All source code, test suites, and benchmark scripts are available in the \texttt{cortex\_v1/} directory of the Livnium repository.\footnote{\url{https://github.com/chetanxpatil/livnium}}

% =========================================================================
\section{Architecture}

The pipeline processes one word at a time through four stages:

\paragraph{Semantic bridge.}
Word $\to$ GloVe-50 vector $\to$ PCA projection into $\mathbb{R}^{3}$ $\to$ unit axis $\hat{n}$ and angle~$\theta$ $\to$ SO(3) rotation matrix $\to$ SU(2) quantum gate.

\paragraph{Lattice.}
The SO(3) rotation is applied to a $3\!\times\!3\!\times\!3$ cubic lattice (27~nodes).  The displacement of each node produces a cosine signal.  The mean absolute cosine across all moved nodes is~$\alpha$.

\paragraph{Governor.}
The polarity governor sets the entropy ceiling at each MPS bond:
\begin{equation}
  S_{\max}^{\mathrm{eff}}(i)
  \;=\;
  S_{\max}^{\mathrm{base}}
  \times
  \bigl(1 + \alpha \times \mathrm{polarity}(i)\bigr).
  \label{eq:governor}
\end{equation}

\paragraph{MPS filter.}
The SU(2) gate is applied to a qubit, followed by a CNOT to an adjacent qubit.  SVD truncation enforces the governor ceiling.  Bonds that exceed the ceiling are pruned; bonds below it survive.

% -------------------------------------------------------------------------
\subsection{The Cubic Lattice}

The lattice consists of 27 nodes at integer coordinates $(x,y,z)$ with $x,y,z\in\{-1,0,1\}$.  Each node carries a symbolic weight $\mathrm{SW}=9f$, where $f$ is the number of coordinates equal to $\pm 1$ (boundary exposure).  The sum $\Sigma\mathrm{SW}=486$ is invariant under all 24~rotations.

$\Sigma\mathrm{SW}$ is a \emph{conserved quantity} of the cubic rotation group: because every rotation is a bijection on the 27 nodes, boundary exposures are permuted but never created or destroyed.  This is the discrete analogue of charge conservation under a continuous symmetry group---the symmetry here being the 24-element subgroup of SO(3) that maps the cube to itself.  The conservation is verified actively: \texttt{is\_valid()} checks $\Sigma\mathrm{SW}=486$ after every rotation and the self-test suite (T4) asserts it across all 24 elements.

The rotation group of the cube has 24 elements distributed across four conjugacy classes (Table~\ref{tab:conjugacy}).  Three generators ($90^{\circ}$ rotations about $x$, $y$, and $z$) produce all 24 elements by BFS, each satisfying $R^{4}=I$.

\begin{table}[h]
  \caption{Conjugacy classes of the cubic rotation group.}
  \label{tab:conjugacy}
  \centering
  \begin{tabular}{lcc}
    \toprule
    Class & Angle & Count \\
    \midrule
    Identity        & $0^{\circ}$   & 1 \\
    Face rotations  & $90^{\circ}$  & 6 \\
    Vertex rotations& $120^{\circ}$ & 8 \\
    Edge rotations  & $180^{\circ}$ & 9 \\
    \bottomrule
  \end{tabular}
\end{table}

% -------------------------------------------------------------------------
\subsection{The Alpha Signal}
\label{sec:alpha}

After a rotation is applied, each node moves from $\mathbf{p}_{\mathrm{old}}$ to $\mathbf{p}_{\mathrm{new}}$.  For every moved node the cosine of the displacement direction relative to the inward radial is
\begin{equation}
  \cos\theta_{i}
  \;=\;
  \frac{(\mathbf{p}_{\mathrm{new}}-\mathbf{p}_{\mathrm{old}})
        \cdot(-\mathbf{p}_{\mathrm{old}})}
       {\lvert\mathbf{p}_{\mathrm{new}}-\mathbf{p}_{\mathrm{old}}\rvert
        \;\lvert\mathbf{p}_{\mathrm{old}}\rvert}\,.
\end{equation}
The $\alpha$ signal is the mean absolute cosine over moved nodes:
\begin{equation}
  \alpha
  \;=\;
  \frac{1}{\lvert\mathcal{M}\rvert}
  \sum_{i\in\mathcal{M}} \lvert\cos\theta_{i}\rvert,
  \qquad
  \mathcal{M}=\{i:\mathbf{p}_{\mathrm{new},i}\neq\mathbf{p}_{\mathrm{old},i}\}.
\end{equation}
Measured $\alpha$ values by class: identity $\alpha=0.000$; face ($90^{\circ}$) $\alpha\approx0.595$; vertex ($120^{\circ}$) $\alpha\approx0.722$; edge ($180^{\circ}$) $\alpha\approx0.836$--$0.841$.

% -------------------------------------------------------------------------
\subsection{The Semantic Bridge}

The bridge maps words to SO(3) rotations using pretrained word embeddings:

\begin{enumerate}[nosep,leftmargin=*]
  \item Load GloVe-50 vectors (50-dimensional; trained on Wikipedia + Gigaword).
  \item Fit PCA on the corpus vocabulary to extract three principal components.
  \item For each word: project the GloVe vector into $\mathbb{R}^{3}$ and normalize to obtain rotation axis $\hat{n}$.
  \item Assign angle~$\theta$ by IDF mode ($\theta=\pi(1-e^{-\mathrm{IDF}(w)})$, rare words get large~$\theta$) or cosine mode ($\theta=\arccos\!\bigl(\cos(v,\bar{v})\bigr)$, words far from the centroid get large~$\theta$).
  \item Construct SO(3) matrix via the axis-angle formula and map to SU(2):
\end{enumerate}
\begin{equation}
  U(\hat{n},\theta)
  \;=\;
  \cos\!\tfrac{\theta}{2}\;\!I
  \;-\;
  i\sin\!\tfrac{\theta}{2}\;(\hat{n}\cdot\boldsymbol{\sigma}),
  \label{eq:su2}
\end{equation}
where $\boldsymbol{\sigma}=(\sigma_{x},\sigma_{y},\sigma_{z})$ are the Pauli matrices.  The mapping is verified homomorphic: $U(A)\,U(B)=U(AB)$ up to global phase, confirmed for 20~random rotation pairs.

A key property: semantically similar words (nearby GloVe vectors) map to nearby rotation axes.  Their SU(2) gates nearly commute, producing constructive interference in the MPS and low entropy growth.  Dissimilar adjacent words produce non-commuting gates and high entropy growth, potentially triggering truncation.

% -------------------------------------------------------------------------
\subsection{The MPS Simulator}

The simulator represents an $n$-qubit state as a tensor train:
\begin{equation}
  \lvert\psi\rangle
  \;=\;
  \sum_{s_{1},\dots,s_{n}}
  A^{[1]}_{s_{1}}\,A^{[2]}_{s_{2}}\cdots A^{[n]}_{s_{n}}\;
  \lvert s_{1}\,s_{2}\cdots s_{n}\rangle,
\end{equation}
where each $A^{[i]}_{s_{i}}$ has shape $(\chi_{\mathrm{left}},\chi_{\mathrm{right}})$ and $s_{i}\in\{0,1\}$.  Memory is $O(n\chi^{2})$ versus $O(2^{n})$ for the full state vector.  When a two-qubit gate increases the bond dimension beyond $\chi_{\max}$, SVD truncation is applied and the discarded weight is logged as truncation error.

Supported operations include Hadamard, Pauli $X/Y/Z$, $R_{x}/R_{y}/R_{z}$, arbitrary single-qubit unitary, CNOT (adjacent and non-adjacent via SWAP chains), and Born-rule measurement with correct left/right environment propagation.

% -------------------------------------------------------------------------
\subsection{The Polarity Governor}

The polarity at bond~$i$ (between qubits $i$ and $i\!+\!1$) in an $n$-qubit system is
\begin{equation}
  \mathrm{polarity}(i)
  \;=\;
  1 - \frac{S(i)}{S_{\max}^{\mathrm{theo}}(i)},
  \qquad
  S_{\max}^{\mathrm{theo}}(i)
  \;=\;
  \min(i+1,\;n-i-1)\times\ln 2,
\end{equation}
where $S(i)$ is the von Neumann entropy at bond~$i$.  The effective ceiling follows Eq.~\eqref{eq:governor}.  High polarity (structured, area-law entanglement) with high~$\alpha$ (semantically important word) yields a generous ceiling; low polarity with low~$\alpha$ yields a tight one.

% =========================================================================
\section{Measured Results: MPS Scaling}
\label{sec:scaling}

\subsection{GHZ State Scaling}

GHZ states ($\lvert 0\rangle^{\otimes n}+\lvert 1\rangle^{\otimes n}$) are area-law: $\chi=2$ regardless of~$n$ (Table~\ref{tab:ghz}).  At $n=1{,}000$, all 100 measurement shots produced valid outcomes (all-0 or all-1).  First-to-last qubit correlation was perfect (200/200 shots: 46\% all-0, 54\% all-1, 0\% mixed).

This is expected MPS behavior.  GHZ is the canonical area-law state and MPS simulators are designed for it.  We report these numbers as a \emph{correctness check}, not as a scaling claim.

\begin{table}[h]
  \caption{GHZ state scaling.  Memory grows linearly; $\chi=2$ at all sizes.}
  \label{tab:ghz}
  \centering
  \begin{tabular}{rccrc}
    \toprule
    $n$ & $\chi_{\max}$ & Time & Memory & Valid \\
    \midrule
       5 & 2 & 0.002\,s &   512\,B & 100/100 \\
      10 & 2 & $<$0.001\,s & 1.2\,KB & 100/100 \\
      50 & 2 & 0.001\,s &  6.3\,KB & 100/100 \\
     100 & 2 & 0.001\,s & 12.7\,KB & 100/100 \\
     500 & 2 & 0.006\,s & 63.9\,KB & 100/100 \\
    1000 & 2 & ---       &127.9\,KB & 100/100 \\
    \bottomrule
  \end{tabular}
\end{table}

% -------------------------------------------------------------------------
\subsection{Random Circuit Scaling}

Random 5-layer brick-wall circuits (alternating even/odd CNOT layers with Hadamard gates) reveal the simulator's limits (Table~\ref{tab:random}).  At $n=20$ the bond dimension hits the $\chi=64$ cap and truncation error accumulates.  Random circuits produce volume-law entanglement that grows exponentially with depth; the MPS approximation degrades accordingly.

\begin{table}[h]
  \caption{Random circuit scaling.  Truncation error appears when $\chi$ saturates.}
  \label{tab:random}
  \centering
  \begin{tabular}{rccc}
    \toprule
    $n$ & $\chi_{\max}$ & Trunc.\ error & Memory \\
    \midrule
    10 & 32         & 0.0  &  43.6\,KB \\
    20 & 64 (cap)   & 10.2 &   1.2\,MB \\
    50 & 64 (cap)   & 10.2 &   5.2\,MB \\
    \bottomrule
  \end{tabular}
\end{table}

% -------------------------------------------------------------------------
\subsection{Governor-Controlled Circuits}

H + CNOT chain + second entangling layer, with governor ($\alpha=0.7$, $S_{\max}=1.5\ln 2$).  Results in Table~\ref{tab:governor}.  Zero truncation error at all sizes.  The governor ceiling is generous enough at high~$\alpha$ and high polarity that no pruning occurs, keeping entanglement in the area-law regime.  This demonstrates that the governor mechanism \emph{functions correctly} at scale; it does not demonstrate that arbitrary circuits can be simulated at 1\,000 qubits.

\begin{table}[h]
  \caption{Governor-controlled circuits.  $\bar{S}$: mean bond entropy (nats); $\bar{P}$: mean polarity.}
  \label{tab:governor}
  \centering
  \begin{tabular}{rccccc}
    \toprule
    $n$ & $\chi$ & Prunes & $\bar{S}$ & $\bar{P}$ & Memory \\
    \midrule
      15 & 4 & 0 & 0.297 & 0.851 &   3.3\,KB \\
     100 & 4 & 0 & 0.336 & 0.967 &  25\,KB \\
     500 & 4 & 0 & 0.344 & 0.990 & 127\,KB \\
    1000 & 4 & 0 & 0.346 & 0.994 & 255\,KB \\
    \bottomrule
  \end{tabular}
\end{table}

% =========================================================================
\section{Measured Results: Organism}

The ChatOrganism processes 1\,360 words through a 15-qubit MPS with $\chi_{\max}=32$ and $S_{\max}=1.5$~bits (Table~\ref{tab:organism}).

The default organism uses MD5 hashing to map words to rotations.  This mapping is deterministic but arbitrary---there is no semantic relationship between a word and its assigned rotation.  Prune/survive decisions are driven by bond entropy thresholds, not by word importance.  The ``survived'' and ``pruned'' labels reflect MPS dynamics, not semantic judgments.

\begin{table}[h]
  \caption{Organism processing statistics (1\,360 words, 15 qubits, $\chi_{\max}\!=\!32$).}
  \label{tab:organism}
  \centering
  \begin{tabular}{lr}
    \toprule
    Metric & Value \\
    \midrule
    Words survived        & 651 (47.9\%) \\
    Words pruned          & 709 (52.1\%) \\
    Governor prune events & 1\,781 \\
    Max $\chi$ used       & 8 \\
    MPS memory            & 8\,896 bytes \\
    Mean bond entropy     & 0.752 nats \\
    Max bond entropy      & 1.152 nats \\
    \bottomrule
  \end{tabular}
\end{table}

% =========================================================================
\section{Measured Results: Alpha-Triage Benchmark}
\label{sec:benchmark}

\subsection{Setup}

A factual paragraph about Shor's algorithm (${\sim}50$~tokens) is streamed through a memory buffer holding only 40\% of the total token count.
Ten ground-truth tokens are identified:
\texttt{1994}, \texttt{peter}, \texttt{shor}, \texttt{algorithm}, \texttt{factor},
\texttt{integers}, \texttt{rsa}, \texttt{encryption}, \texttt{decoherence}, \texttt{qubits}.
Three eviction policies compete under identical capacity:

\begin{itemize}[nosep,leftmargin=*]
  \item \textbf{FIFO}: evict the oldest token.
  \item \textbf{LRU}: evict the least-recently-used token (equivalent to FIFO in a single-pass stream with no re-reads).
  \item \textbf{Alpha-triage}: evict the token with the lowest $\alpha$ value.
\end{itemize}

% -------------------------------------------------------------------------
\subsection{Results}

\begin{table}[h]
  \caption{Fact recall under three eviction policies (single Shor-paragraph benchmark).}
  \label{tab:benchmark}
  \centering
  \begin{tabular}{lc}
    \toprule
    Policy & Fact Recall \\
    \midrule
    FIFO                    & 30\% (3/10) \\
    LRU                     & 30\% (3/10) \\
    $\alpha$-triage (mock)  & 100\% (10/10) \\
    \bottomrule
  \end{tabular}
\end{table}

In mock mode, $\alpha$ values are assigned by a hand-crafted lookup table: high-information content words receive $\alpha=0.95$, mid-range words $\alpha=0.50$, common function words $\alpha=0.25$.  In real mode (GloVe-derived $\alpha$ via the semantic bridge), the same qualitative result holds with noisier margins.

FIFO and LRU retain only the three ground-truth facts that happen to appear in the last 40\% of the passage.  $\alpha$-triage retains all ten because factual tokens have high~$\alpha$ (they are rare, semantically distinctive, and far from the GloVe corpus centroid) and are therefore evicted last.

% =========================================================================
\section{Measured Results: TFIM Quantum Quench Benchmark}
\label{sec:quench}

\paragraph{Setup.}
We benchmark the polarity governor on a concrete quantum dynamics task: real-time evolution under the transverse-field Ising model (TFIM),
\begin{equation}
  H = -J \sum_{i} Z_i Z_{i+1} \;-\; h \sum_{i} X_i,
\label{eq:tfim}
\end{equation}
with $J=1.0$, $h=0.5$, $n=14$ qubits, evolved for $T=1.0$ in 10 first-order Trotter steps ($dt=0.1$).
Initial state: $|\!\uparrow\cdots\uparrow\rangle$ ($Z$-polarized product state).
Three backends compete under identical $\chi_{\max}=8$ and $S_{\text{base}}=0.70$~nats: a reference simulation ($\chi=64$, no truncation), a fixed-entropy governor ($\alpha=0$), and the polarity-aware governor ($\alpha=0.8$).
Observables are exact $\langle Z_i\rangle$ computed via left/right environment contractions (no shot noise).

\paragraph{Results.}
Table~\ref{tab:quench} reports mean absolute error (MAE) in $\langle Z_i\rangle$ versus the reference at $t=1.0$.

\begin{table}[h]
  \caption{TFIM quench benchmark. MAE is mean $|\langle Z_i\rangle_{\text{backend}} - \langle Z_i\rangle_{\text{ref}}|$ over all 14 sites.}
  \label{tab:quench}
  \centering
  \small
  \begin{tabular}{lccc}
    \toprule
    Backend & MAE vs.\ ref & Prune events & $\chi_{\max}$ reached \\
    \midrule
    Reference ($\chi=64$) & 0.0000 & — & 64 \\
    Fixed Governor        & 0.2193 & 39 & 6  \\
    Polarity Governor     & 0.1245 & 72 & 8  \\
    \bottomrule
  \end{tabular}
\end{table}

The polarity governor reduces MAE by \textbf{43.2\%} relative to the fixed ceiling, and 8 of 14 sites are closer to the reference.
The fixed governor makes 39 coarse prune decisions; the polarity governor makes 72 finer decisions and allows high-polarity bonds to expand toward $\chi=8$, preserving structured entanglement in the volume-law core of the chain while aggressively truncating area-law boundary bonds.

\paragraph{Physical interpretation.}
The $\alpha$ signal derived from the polarity score biases truncation toward volume-law bonds, protecting the magnetization profile in the dynamically entangled interior (sites 3--12) while pruning the near-product boundary sites.
This is the predicted behavior of the governor: geometry-derived importance determines \emph{which} bonds absorb the entropy budget, not merely \emph{how much} budget is available.

% -------------------------------------------------------------------------
\subsection{Limitations of the Benchmark}

We state these explicitly:

\begin{enumerate}[nosep,leftmargin=*]
  \item The mock $\alpha$ table is hand-crafted.  That hand-labelled importance beats content-blind eviction is expected.
  \item The passage is short ($\sim$50 tokens) and facts are front-loaded, which disadvantages FIFO/LRU.
  \item The ground truth is hand-selected; different fact sets yield different recall numbers.
  \item The experiment is single-pass.  In a multi-access scenario, LRU would outperform FIFO.  $\alpha$-triage would still outperform LRU because importance and recency are different signals.
  \item One paragraph, one topic, one ground truth.  This is a proof-of-concept, not a corpus-level evaluation.
\end{enumerate}

% =========================================================================
\section{Ramsey R(3,3) Geometric Sandbox}
\label{sec:ramsey}

\paragraph{Motivation.}
To validate that the cubic lattice tension signal transfers from quantum entropy to
combinatorial constraint satisfaction, we map the complete graph $K_6$ onto the
$3\!\times\!3\!\times\!3$ Livnium lattice and test whether a purely geometric pressure
field can detect and localize forbidden monochromatic triangles ($K_3$).

\paragraph{Vertex placement.}
Six vertices of $K_6$ are placed at the six face-center cells of the cube
(the unique octahedral arrangement where all vertices carry identical symbolic
weight $\mathrm{SW}=9$ and no topological distortion is introduced).
The 15 edges map to 12 distinct edge-type cells plus the center cell for the 3 axis-aligned pairs.

\paragraph{R(3,3) detection (discrete injection).}
Tension is injected directly into the three mediating cells of any monochromatic triangle
(+10 per cell; no Gaussian blur, no triangle-check code).
Tested across all $2^{15}=32{,}768$ 2-colorings of $K_6$: \textbf{32,768/32,768 pass} —
the lattice peak always falls inside a monochromatic triangle's edge-cell set.
An earlier continuous Gaussian field ($\sigma^2=0.5$) gave 0.05\% accuracy due to all
12 interior triangle centroids aliasing into the center cell; switching to strict discrete
injection eliminated all false positives.

\paragraph{Self-healing lattice (emergent nonlinear tension).}
We remove all explicit triangle checks and assign each edge cell a charge
($\mathrm{Red}=+1$, $\mathrm{Blue}=-1$).  Tension at any cell is the square of its
von-Neumann neighborhood sum.  A monochromatic $K_3$ produces a local sum of $\pm3$
($T=9$); a mixed triangle produces $\pm1$ ($T=1$).  The field ``sees'' the forbidden
structure without being told to look.

A Simulated Annealing healer flips edges to minimize total tension, accepting worsening
moves with probability $\exp(-\Delta T/\tau)$ to escape local minima.
Results over 100/200 random restarts:

\begin{table}[h]
  \caption{Self-healing Ramsey lattice results. Phase A: $K_5$ (valid colorings exist).
           Phase B: $K_6$ on $3\!\times\!3\!\times\!3$ lattice (R(3,3)=6, no valid coloring).}
  \label{tab:ramsey}
  \centering
  \small
  \begin{tabular}{llccc}
    \toprule
    Phase & Method & Valid finds & $T_{\min}$ & mono-$K_3$ min \\
    \midrule
    A ($K_5$) & Greedy & 39/100 (39\%)  & —      & 0 \\
    A ($K_5$) & SA     & 100/100 (100\%) & —     & 0 \\
    B ($K_6$) & Greedy & 0 T=0           & 27.0   & 2 \\
    B ($K_6$) & SA     & 0 T=0           & $27.0\pm 0.0$ & 2 \\
    \bottomrule
  \end{tabular}
\end{table}

The SA healer finds a valid $K_5$ coloring on every random start and never finds $T=0$
for $K_6$ (consistent with $R(3,3)=6$).
The $K_6$ lattice floor is $T=27.0$ with zero variance — the SA reliably locates the
global minimum of the spatial energy landscape.

\paragraph{Interpretation and limitations.}
The system is \emph{complete} (finds solutions when they exist) and \emph{sound}
(does not hallucinate solutions where none exist).
However, the lattice tension minimum ($T=27$) and the combinatorial minimum
(mono-$K_3 = 2$) are not identical objectives: some colorings with 2 monochromatic
triangles carry higher lattice tension due to neighborhood-sum interference from
spatially dispersed triangle edges.
The field is most sensitive to compact triangles (spatially adjacent edge-cells) and
underestimates tension from triangles whose edges span distant lattice regions.
This is an honest limitation of projecting a 15-edge graph onto 27 cells.

% =========================================================================
\section{R(5,5) Heuristic Search — Triangle Adjacency Kernel}
\label{sec:r55}

\paragraph{Motivation.}
Following the R(3,3) sandbox, we extend to the open problem $R(5,5)$, whose exact value is
unknown (current bounds: $43 \le R(5,5) \le 48$~\cite{exoo1989,mckay1995}).
A certificate that $R(5,5)>n$ requires an explicit 2-coloring of $K_n$ with no monochromatic $K_5$.
We do not claim to find such a certificate.
Instead, we benchmark three heuristic tension kernels for guiding simulated annealing (SA)
search toward low mono-$K_5$ colorings, as a test of whether the geometric signal quality
transfers from the $R(3,3)$ sandbox to a harder open problem.

\paragraph{Kernel evolution.}
We designed and tested three tension functions $T$ over edge colorings of $K_n$:

\begin{enumerate}[nosep,leftmargin=*]
\item \textbf{Vertex-sum kernel} $T=\sum_v \sigma(v)^2$ (where $\sigma(v)=\deg_R(v)-\deg_B(v)$):
  SA immediately found $T=0$ on $K_{45}$, appearing to prove $R(5,5)>45$.
  This is a \emph{parity loophole}: $T=0$ requires only balanced degree per vertex, which is
  achievable on any even-degree graph regardless of clique content.  Rejected.

\item \textbf{Cubic density kernel} $T=\sum_v [\sigma_R(v)^3 + \sigma_B(v)^3]$:
  Eliminates the parity loophole.  Phase A validation passes ($K_3/K_4/K_5$ all correct).
  Phase B: all 20 restarts on $K_{43}$ converge to mono-$K_5\approx 736$, with zero variance.
  Diagnosis: a $K_5$ contributes only 4 out of 42 vertex edges ($\approx9.5\%$),
  making the $K_5$ signal-to-background ratio $\approx1\!:\!10$.

\item \textbf{Triangle adjacency kernel} $T=\sum_{(u,v)}[t_R(u,v)^3 + t_B(u,v)^3]$,
  where $t_R(u,v)=|\{w : e(u,w)=R,\,e(v,w)=R\}|$ (common red neighbours):
  A $K_5$ edge has $t_R\ge3$, yielding a cubic spike $\ge27$ versus $\approx1$--$2$ for a
  typical edge.  Signal fraction $\approx27\%$ above baseline per $K_5$ edge, versus
  $\approx10\%$ for the cubic kernel.  Update cost $O(n)$ per flip via incremental
  maintenance of the $t_R,t_B$ tables (built once via $A_R \cdot A_R$ matrix multiplication).
\end{enumerate}

\paragraph{Results.}
Phase A validation confirms that the triangle kernel correctly identifies the $R(3,3)=6$
barrier: $K_5$ finds 8/20 (40\%) $K_3$-free colorings (valid; $R(3,3)>5$) and $K_6$ finds 0/20
(correct; $R(3,3)=6$).
The 40\% success rate on $K_5$ (vs.\ 100\% for the cubic kernel) indicates the SA temperature
schedule requires retuning for the steeper triangle-kernel landscape; this is a known limitation.

Phase B results on 10 independent SA restarts per graph are shown in Table~\ref{tab:r55}.

\begin{table}[h]
  \caption{R(5,5) heuristic search: triangle adjacency kernel vs.\ cubic kernel.
           Lower mono-$K_5$ is better; 0 would be a certificate.
           Reported as best across 10 independent SA restarts.
           \textbf{These are heuristic results, not proofs.}}
  \label{tab:r55}
  \centering
  \small
  \begin{tabular}{lcccc}
    \toprule
    Graph & Total $K_5$ subgraphs & Cubic best & Triangle best & Improvement \\
    \midrule
    $K_{43}$ & 962,598   & $\approx736$ & \textbf{484} & 34\% \\
    $K_{44}$ & 1,086,008 & $\approx875$ & \textbf{573} & 35\% \\
    \bottomrule
  \end{tabular}
\end{table}

The triangle kernel achieves a consistent $\approx34$--$35\%$ reduction in monochromatic $K_5$
count relative to the cubic kernel, with the improvement holding across both graph sizes.
The fraction of monochromatic $K_5$ subgraphs remaining is $\approx0.050$--$0.053\%$
of all $K_5$ subgraphs, stable across $K_{43}$ and $K_{44}$.
SA restarts show genuine variance in mono-$K_5$ (9--11\% spread), confirming the search
is exploring distinct local minima rather than collapsing to a single basin.

\paragraph{Interpretation and limitations.}
\emph{No result here constitutes a proof about $R(5,5)$.}
A proof would require mono-$K_5=0$, which we do not achieve.
The literature lower bound ($R(5,5)\ge43$) was established by Exoo (1989) via explicit
certificate; we do not reproduce or improve it.
The contribution is a \emph{kernel comparison}: the 2-hop triangle adjacency signal is a
more effective heuristic guide for SA search than the 1-hop cubic density signal,
reducing the monochromatic clique floor by $\approx34\%$ on the graphs tested.
Future directions include longer SA schedules, population-based search with coloring
cross-pollination, and targeted edge reheating around highest-tension edges.

% =========================================================================
\section{Verified Invariants}

The self-test suite verifies 12 structural invariants; all 12 pass:

\begin{enumerate}[nosep,leftmargin=*]
  \item $R^{4}=I$ for all three generators.
  \item 24 distinct orientations generated by BFS.
  \item Class counts: $\{0\!:\!1,\;1\!:\!6,\;2\!:\!12,\;3\!:\!8\}$.
  \item $\Sigma\mathrm{SW}=486$ invariant under all 24 rotations.
  \item Lattice bijection preserved after every rotation.
  \item $\cos\theta\in[-1,+1]$ for all polarity signals.
  \item Identity rotation $\Rightarrow$ $\alpha=0$.
  \item $\alpha$ range $[0.595,\,0.841]$, std $=0.094$.
  \item High-$\alpha$ governor prunes $\leq$ low-$\alpha$ governor prunes.
  \item Homomorphism: $U(A)\,U(B)\approx U(AB)$ for 20 random pairs.
  \item All 24 SU(2) gates unitary and distinct.
  \item GHZ-15: 200/200 valid measurement outcomes.
\end{enumerate}

% =========================================================================
\section{What This System Does Not Do}

\begin{enumerate}[nosep,leftmargin=*]
  \item The MPS state is a quantum state.  It does not store words or retrieve them.  It is a compressed encoding of the sequence of gates applied.
  \item The organism's prune/survive decisions are driven by bond entropy relative to the governor ceiling, not by semantic understanding.
  \item The MD5 $\to$ rotation mapping (used in the default organism) is arbitrary.  Only the semantic bridge (GloVe mode) connects word meaning to lattice geometry.
  \item GHZ scaling to 1\,000 qubits is standard MPS behavior for area-law states.  It is not a demonstration of general large-scale quantum simulation.
  \item The polarity governor adjusts entropy ceilings but cannot increase the hard bond dimension cap~$\chi_{\max}$.
  \item The $\alpha$-triage benchmark is a single-paragraph, single-topic, hand-labelled experiment.  Generalization to other domains is untested.
\end{enumerate}

% =========================================================================
\section{Related Work}

\paragraph{MPS simulators.}
ITensor~\citep{itensor}, TeNPy~\citep{tenpy}, and Quimb~\citep{quimb} are mature tensor network libraries.  They provide standard MPS operations (SVD truncation, DMRG, time evolution) but do not include content-aware governor mechanisms.

\paragraph{Cache eviction.}
FIFO, LRU, LFU, ARC~\citep{arc}, and LIRS~\citep{lirs} are well-studied eviction policies.  Content-aware caching has been explored in the CDN literature.  The $\alpha$-triage policy differs in that the importance signal is derived from a geometric rotation group rather than from access patterns or file sizes.

\paragraph{Word embeddings for information retrieval.}
TF-IDF, BM25, and dense retrieval methods (DPR, ColBERT) use embedding similarity to rank documents.  The semantic bridge uses embeddings differently: projecting into SO(3) to produce a scalar signal that controls a quantum simulator's entropy budget.

% =========================================================================
\section{Conclusion}

We presented a system that couples a cubic lattice to an MPS simulator via a polarity governor.  The measured results are: (i)~1\,000-qubit GHZ states with perfect measurement validity (area-law, $\chi=2$); (ii)~zero truncation error at 1\,000 qubits under governor control ($\chi=4$, polarity $>0.99$); (iii)~100\% fact recall under $\alpha$-triage versus 30\% for FIFO/LRU on a single benchmark paragraph; (iv)~43.2\% reduction in TFIM magnetization error under the polarity governor versus a fixed entropy ceiling at identical $\chi_{\max}$; (v)~a Ramsey R(3,3) sandbox: 32,768/32,768 $K_3$ violations detected, SA healer 100\% on $K_5$, confirmed floor above $T=0$ for $K_6$; (vi)~an R(5,5) heuristic benchmark: triangle adjacency kernel achieves best mono-$K_5=484$ on $K_{43}$ and $573$ on $K_{44}$, a 34--35\% improvement over the cubic density kernel (no certificate claimed); and (vii)~12/12 verified invariants.

Whether these results constitute a useful contribution is for the reader to assess.  The code is open, the tests are runnable, and the numbers are reproducible.  We have tried to state what was measured without overstating what it means.

% =========================================================================
\bibliographystyle{plainnat}
\begin{thebibliography}{5}

\bibitem[Fishman et~al.(2022)]{itensor}
M.~Fishman, S.~R.~White, and E.~M.~Stoudenmire.
\newblock The {IT}ensor software library for tensor network calculations.
\newblock \emph{SciPost Phys.\ Codebases}, 4, 2022.

\bibitem[Hauschild \& Pollmann(2018)]{tenpy}
J.~Hauschild and F.~Pollmann.
\newblock Efficient numerical simulations with tensor networks: {T}ensor {N}etwork {P}ython ({TeNPy}).
\newblock \emph{SciPost Phys.\ Lect.\ Notes}, 5, 2018.

\bibitem[Gray(2018)]{quimb}
J.~Gray.
\newblock quimb: A {P}ython library for quantum information and many-body calculations.
\newblock \emph{J.\ Open Source Softw.}, 3(29):819, 2018.

\bibitem[Megiddo \& Modha(2003)]{arc}
N.~Megiddo and D.~S.~Modha.
\newblock {ARC}: A self-tuning, low overhead replacement cache.
\newblock In \emph{FAST}, 2003.

\bibitem[Jiang \& Zhang(2002)]{lirs}
S.~Jiang and X.~Zhang.
\newblock {LIRS}: An efficient low inter-reference recency set replacement policy.
\newblock In \emph{ACM SIGMETRICS}, 2002.

\end{thebibliography}

\end{document}
